\chapter[1949 --- Moniz]{1949: Walter Rudolf Hess and Ant\'onio Egas Moniz}
\label{ch:1949_Hess}

\section*{Main Contribution}

\textit{Discovered the coordination function of the midbrain reticular formation, explaining how the brain regulates wakefulness and arousal.}

\section{Official Citation}

for his discovery of the functional organization of the interbrain as a coordinator of the activities of the internal organs; and for his discovery of the therapeutic value of leucotomy in certain psychoses

\section{Background and Historical Context}

\subsection{Walter Rudolf Hess}

Walter Rudolf Hess was born on March 17, 1881, in Frauenfeld, Switzerland. He grew up in a family that valued education and public service; his father was a physician, and young Walter followed in his footsteps, enrolling at the University of Zurich Medical School in 1900. Hess received his medical degree in 1906 and subsequently trained in surgery and internal medicine at several Swiss hospitals. However, his interest soon turned to physiology, and he joined the Physiological Institute at the University of Zurich, where he would spend the next several decades conducting research on the nervous system.

Hess's research focused on the diencephalon---the region of the brain that includes the thalamus and hypothalamus---and its role in the regulation of autonomic functions. Working with cats as his experimental model, he developed a technique for electrical stimulation of specific brain regions using precisely placed electrodes, a method that allowed him to map the functional organization of the diencephalon with notable precision.

The scientific context of the early twentieth century was one of growing interest in the relationship between the brain and the body's internal organs. The autonomic nervous system---the division of the nervous system that controls involuntary functions such as heart rate, blood pressure, digestion, and temperature regulation---had been described by John Newport Langley in the early 1900s, but the central mechanisms that coordinate autonomic function remained poorly understood. It was within this intellectual framework that Hess conducted his pioneering research.

\subsection{Ant\'onio Egas Moniz}

Ant\'onio Caetano de Abreu Freire Egas Moniz was born on November 29, 1874, in {\AA}vora Monte, near Lous{\~a}, Portugal. He came from a wealthy landowning family and received his education in medicine at the University of Coimbra, where he graduated with honors in 1899. Moniz initially pursued a career in neurology, making important contributions to the diagnosis of brain tumors using contrast radiography. He invented the technique of cerebral angiography---the injection of a radiopaque contrast agent into the blood vessels of the brain to visualize them on X-ray---which became an essential diagnostic tool in neurology.

Moniz also had a parallel career in politics, serving as Portugal's Minister of Foreign Affairs and playing a key role in Portuguese diplomacy during and after World War I. He returned to medicine full-time in the 1920s and became interested in the treatment of severe psychiatric disorders, which at the time had very few effective therapies. Patients with schizophrenia, severe depression, and other psychotic disorders were often confined to asylums for life, with little hope of recovery.

Moniz was influenced by the work of several researchers who had observed that brain damage or disease could sometimes lead to improvements in psychiatric symptoms. He hypothesized that by intentionally severing certain neural connections in the brain, it might be possible to alleviate the symptoms of severe mental illness. This hypothesis led him to develop the surgical procedure known as leucotomy (later called lobotomy).

\section{The Discovery/Contribution}

The Nobel Prize was awarded to Hess and Moniz in different areas of neuroscience. Hess received the prize ``for his discovery of the functional organization of the interbrain as a coordinator of the activities of the internal organs,'' while Moniz received it ``for his discovery of the therapeutic value of leucotomy in certain psychoses.'' While these contributions are quite different in nature, both address fundamental questions about the relationship between the brain and bodily function.

\subsection{Hess's Mapping of the Diencephalon}

Hess's primary contribution was the demonstration that the diencephalon contains a highly organized system for the coordination of autonomic functions. Using his electrical stimulation technique, Hess systematically explored the effects of stimulating different regions of the diencephalon and produced a detailed functional map of this brain region.

He showed that stimulation of specific areas of the hypothalamus could elicit a wide range of autonomic responses, including changes in heart rate, blood pressure, body temperature, respiration, and digestive function. For example, stimulation of the anterior hypothalamus caused cooling of the body (by promoting sweating and vasodilation), while stimulation of the posterior hypothalamus caused heating (by promoting vasoconstriction and shivering). Stimulation of certain hypothalamic areas could also produce complex behavioral responses, such as the ``defense reaction'' (increased heart rate, elevated blood pressure, and piloerection) or the ``feeding response'' (searching for and consuming food).

Hess's experiments demonstrated that the hypothalamus serves as the master regulator of the body's internal environment, coordinating the activity of the autonomic nervous system and the endocrine system to maintain homeostasis. He also showed that the diencephalon is organized in a hierarchical fashion, with higher brain regions (such as the cerebral cortex) exerting regulatory control over the autonomic functions coordinated by the hypothalamus.

One of a notable aspects of Hess's work was his discovery of the ``sleep-wakefulness'' cycle controlled by the diencephalon. He showed that stimulation of certain hypothalamic areas could induce sleep-like states, while stimulation of other areas produced wakefulness. These observations anticipated by several decades the discovery of the neural circuits that regulate sleep and wakefulness, which is now known to involve specific populations of neurons in the hypothalamus and brainstem.

\subsection{Moniz and the Leucotomy}

Moniz's contribution was the development of leucotomy (lobotomy), a surgical procedure in which the white matter connections between the prefrontal cortex and other brain regions are severed. Moniz performed the first leucotomy in 1935, in collaboration with the neurosurgeon Almeida Lima, on a patient with severe schizophrenia. The procedure involved injecting alcohol into the frontal lobes to destroy the neural connections between the prefrontal cortex and the rest of the brain.

The results were, by the standards of the time, remarkable. Several of Moniz's patients, who had been severely incapacitated by their psychiatric symptoms, showed significant improvement after the procedure. They became calmer, less agitated, and more manageable, and some were able to leave the asylum and return to their families. Moniz reported his results in a series of publications and argued that leucotomy was a valuable therapeutic option for patients with severe mental illness who had not responded to other treatments.

Moniz's work was both celebrated and controversial. His supporters argued that leucotomy offered hope to patients who would otherwise be condemned to a lifetime of institutional care. His critics argued that the procedure was too radical, that its effects were unpredictable, and that it often resulted in significant cognitive and personality changes. The controversy surrounding leucotomy would intensify in the 1940s and 1950s, when the American neurologist Walter Freeman popularized the transorbital lobotomy---a simplified version of the procedure that could be performed quickly and without specialized surgical equipment---and performed thousands of operations on patients across the United States.

\section{Impact on Modern Medicine}

The contributions of Hess and Moniz have had very different impacts on modern medicine, reflecting the very different natures of their discoveries.

\subsection{Hess's Contributions to Neurophysiology and Neurology}

Hess's mapping of the diencephalon has been fundamental to our understanding of the regulation of autonomic functions and the pathophysiology of neurological and psychiatric disorders. His demonstration that the hypothalamus coordinates the body's internal environment has been essential for the diagnosis and treatment of a wide range of conditions, including hypothalamic disorders, sleep disorders, and autonomic dysfunction.

Understanding the hypothalamic regulation of body temperature has been crucial for the management of fever and hypothermia. The recognition that the hypothalamus controls the circadian rhythm of sleep and wakefulness has led to the development of treatments for sleep disorders, including light therapy for seasonal affective disorder and melatonin supplementation for jet lag.

Hess's work on the hypothalamic regulation of endocrine function has also had important implications for endocrinology. The hypothalamus controls the pituitary gland through the release of hypothalamic hormones, and dysfunction of this regulatory system can lead to conditions such as diabetes insipidus, hypothyroidism, and growth hormone deficiency. Understanding the hypothalamic-pituitary axis has been essential for the development of hormone replacement therapies for these conditions.

\subsection{Moniz and the Rise and Fall of Psychosurgery}

The impact of Moniz's leucotomy on modern medicine is more complex and controversial. On the one hand, the development of leucotomy represented the first attempt to treat severe psychiatric disorders by directly modifying brain structure. This approach---now known as psychosurgery---opened up an entirely new therapeutic frontier and demonstrated that the brain could be surgically modified to alleviate psychiatric symptoms.

On the other hand, the widespread and often indiscriminate use of leucotomy and lobotomy in the 1940s and 1950s caused significant harm to many patients. It is estimated that approximately 40,000 lobotomies were performed in the United States alone during this period, and many patients suffered severe cognitive and personality changes as a result. The introduction of chlorpromazine and other antipsychotic drugs in the 1950s provided a pharmacological alternative to surgery and led to a dramatic decline in the use of lobotomy.

The legacy of Moniz's leucotomy is also complicated by ethical concerns. Many of the patients who underwent lobotomy were unable to provide informed consent, as they were severely incapacitated by their psychiatric illness. The use of lobotomy on involuntary patients, including children and prisoners, has been widely condemned as a violation of human rights.

Despite these controversies, Moniz's work has had a lasting impact on the field of psychosurgery. Modern techniques such as deep brain stimulation (DBS) and capsulotomy---which involve targeted modification of specific brain circuits---represent more refined and ethically rigorous versions of the approach that Moniz pioneered. These techniques are currently being investigated for the treatment of severe, treatment-resistant depression, obsessive-compulsive disorder, and other psychiatric conditions.

\section{Legacy}

The legacies of Walter Hess and Ant\'onio Egas Moniz reflect the diverse and sometimes contradictory nature of scientific progress. Hess's meticulous mapping of the diencephalon exemplifies the power of careful, systematic experimentation to reveal the inner workings of the nervous system. His work has had an enduring impact on neurophysiology, neurology, and endocrinology, and his techniques have been adopted and refined by generations of neuroscientists.

Hess continued his research at the University of Zurich until his retirement and received numerous honors during his long career. He died on August 12, 1973, at the age of 92.

Moniz's legacy is more complex. His development of leucotomy was both a pioneering achievement and a cautionary tale. The recognition that brain surgery could alleviate psychiatric symptoms was an important advance, but the widespread and often indiscriminate use of the procedure caused significant harm. Moniz was shot and seriously wounded by a former patient in 1939, an event that underscored the controversial nature of his work. He continued to advocate for leucotomy until his death on December 13, 1955, at the age of 81.

The stories of Hess and Moniz illustrate the challenges and responsibilities that come with scientific and medical innovation. Their work has contributed to our understanding of the brain and its disorders, but it has also raised important questions about the ethics of medical experimentation, the limits of intervention in the nervous system, and the need for rigorous evaluation of new therapies before they are widely adopted.


\section{Modern Developments}

Hess and Moniz's work has evolved into modern neuropsychiatry. Understanding of the reticular formation has informed treatments for disorders of consciousness (coma, minimally conscious state). Deep brain stimulation of the thalamus is being explored for consciousness disorders. Understanding of arousal circuits has led to treatments for narcolepsy (orexin receptor agonists) and insomnia (duorexant). Transcranial magnetic stimulation (TMS) modulates cortical excitability for depression and other neuropsychiatric conditions.
